\documentclass[11pt]{article}

% Change "review" to "final" to generate the final (sometimes called camera-ready) version.
% Change to "preprint" to generate a non-anonymous version with page numbers.
\usepackage[preprint]{acl}

% Standard package includes
\usepackage{times}
\usepackage{latexsym}

% For proper rendering and hyphenation of words containing Latin characters (including in bib files)
\usepackage[T1]{fontenc}
% For Vietnamese characters
% \usepackage[T5]{fontenc}
% See https://www.latex-project.org/help/documentation/encguide.pdf for other character sets

% This assumes your files are encoded as UTF8
\usepackage[utf8]{inputenc}

% This is not strictly necessary, and may be commented out,
% but it will improve the layout of the manuscript,
% and will typically save some space.
\usepackage{microtype}

% This is also not strictly necessary, and may be commented out.
% However, it will improve the aesthetics of text in
% the typewriter font.
\usepackage{inconsolata}
\usepackage{amsmath}

%Including images in your LaTeX document requires adding
%additional package(s)
\usepackage{graphicx}

% If the title and author information does not fit in the area allocated, uncomment the following
%
%\setlength\titlebox{<dim>}
%
% and set <dim> to something 5cm or larger.

\title{Query-Adaptive Learned Fusion for Robust Multimodal Retrieval-Augmented Generation}

% Author information can be set in various styles:
% For several authors from the same institution:
% \author{Author 1 \and ... \and Author n \\
%         Address line \\ ... \\ Address line}
% if the names do not fit well on one line use
%         Author 1 \\ {\bf Author 2} \\ ... \\ {\bf Author n} \\
% For authors from different institutions:
% \author{Author 1 \\ Address line \\  ... \\ Address line
%         \And  ... \And
%         Author n \\ Address line \\ ... \\ Address line}
% To start a separate ``row'' of authors use \AND, as in
% \author{Author 1 \\ Address line \\  ... \\ Address line
%         \AND
%         Author 2 \\ Address line \\ ... \\ Address line \And
%         Author 3 \\ Address line \\ ... \\ Address line}

\author{Ishani Kathuria \\
  College of Technology \\
  Purdue University Northwest \\
  Hammond, Indiana, USA \\
  \texttt{ishani@kathuria.net} \\\And
  Manghui Tu \\
  College of Technology \\
  Purdue University Northwest \\
  Hammond, Indiana, USA \\
  \texttt{tu.15@pnw.edu} \\}

%\author{
%  \textbf{First Author\textsuperscript{1}},
%  \textbf{Second Author\textsuperscript{1,2}},
%  \textbf{Third T. Author\textsuperscript{1}},
%  \textbf{Fourth Author\textsuperscript{1}},
%\\
%  \textbf{Fifth Author\textsuperscript{1,2}},
%  \textbf{Sixth Author\textsuperscript{1}},
%  \textbf{Seventh Author\textsuperscript{1}},
%  \textbf{Eighth Author \textsuperscript{1,2,3,4}},
%\\
%  \textbf{Ninth Author\textsuperscript{1}},
%  \textbf{Tenth Author\textsuperscript{1}},
%  \textbf{Eleventh E. Author\textsuperscript{1,2,3,4,5}},
%  \textbf{Twelfth Author\textsuperscript{1}},
%\\
%  \textbf{Thirteenth Author\textsuperscript{3}},
%  \textbf{Fourteenth F. Author\textsuperscript{2,4}},
%  \textbf{Fifteenth Author\textsuperscript{1}},
%  \textbf{Sixteenth Author\textsuperscript{1}},
%\\
%  \textbf{Seventeenth S. Author\textsuperscript{4,5}},
%  \textbf{Eighteenth Author\textsuperscript{3,4}},
%  \textbf{Nineteenth N. Author\textsuperscript{2,5}},
%  \textbf{Twentieth Author\textsuperscript{1}}
%\\
%\\
%  \textsuperscript{1}Affiliation 1,
%  \textsuperscript{2}Affiliation 2,
%  \textsuperscript{3}Affiliation 3,
%  \textsuperscript{4}Affiliation 4,
%  \textsuperscript{5}Affiliation 5
%\\
%  \small{
%    \textbf{Correspondence:} \href{mailto:email@domain}{email@domain}
%  }
%}

\begin{document}
\maketitle
\begin{abstract}
This document is a supplement to the general instructions for *ACL authors. It contains instructions for using the \LaTeX{} style files for ACL conferences.
The document itself conforms to its own specifications, and is therefore an example of what your manuscript should look like.
These instructions should be used both for papers submitted for review and for final versions of accepted papers.
\end{abstract}

\section{Introduction}

Retrieval-augmented generation (RAG) has emerged as a powerful paradigm to ground large language model (LLM) outputs in external knowledge, significantly reducing hallucination and enabling factually accurate responses across diverse domains. While initial RAG systems relied on single retrieval modalities, either dense vector search for semantic similarity or sparse keyword indexing for lexical precision, recent work demonstrates that hybrid retrieval, combining multiple orthogonal ranking signals, achieves substantially better recall, robustness, and adaptability \cite{ahmad_benchmarking_2025}. Systems like HybGRAG \cite{lee_hybgrag_2025} and RAG-Fusion \cite{rackauckas_rag-fusion_2024} have shown that multi-source fusion outperforms single modalities, establishing hybrid retrieval as a best practice for modern RAG systems.

However, contemporary hybrid RAG systems face a critical limitation: they employ static, uniform weighting schemes across all queries. Whether using fixed fusion parameters or query decomposition with Reciprocal Rank Fusion (RRF), existing approaches fail to adapt to the heterogeneous nature of real-world information needs. A straightforward factual lookup ("Who won the 2024 World Cup?") should prioritize different retrieval modalities than a complex comparative question ("Compare GPU architectures from 2022 to 2024 in terms of performance and energy efficiency"). Yet current systems treat both queries identically. Worse, when documents are multimodal that is they contain text, tables, and images, existing fusion mechanisms lack explicit mechanisms to validate cross-modal consistency or detect suspicious evidence patterns.

An emerging but underexplored concern compounds these limitations: adversarial robustness. MM-PoisonRAG \cite{ha_mm-poisonrag_2025} recently demonstrated that multimodal RAG systems are uniquely vulnerable to knowledge poisoning attacks, achieving 56\% success rates even with restricted attacker access. The paper shows that cross-modal alignment can be systematically exploited with attackers crafting image-caption pairs that rank highly in similarity scores but contain false information. The robustness evaluation studies \cite{amirshahi_evaluating_2025} make a crucial observation that diverse evidence sources provide natural defense against such attacks. Critically, existing defenses like adaptive gating, layer-knowledge attention \cite{shi_making_2025}, and post-hoc conflict resolution operate at the generation stage, which is after poisoned evidence has already been retrieved and ranked. This creates a fundamental asymmetry: adversaries attack at the retrieval-fusion stage, but defenders react at the generation stage.

QALF, Query-Adaptive Learned Fusion, addresses these challenges by introducing a foundational architecture that simultaneously achieves three goals:
\begin{enumerate}
    \item \textbf{Performance Optimization}: Dynamically adapts retrieval fusion weights based on query properties (complexity, intent and modality requirements), extending prior work on query-adaptive routing \cite{jeong_adaptive-rag_2024} from generation-stage strategy selection to retrieval-stage weight assignment. This enables simple factual queries to leverage efficient, vector-dominant retrieval while complex multi-hop questions benefit from graph traversal and keyword matching.
    \item \textbf{Multimodal Robustness}: Implements consensus-based validation where documents are weighted by agreement across independent retrieval modalities. Because poisoned evidence typically exploits weaknesses in a single retriever, QALF's consensus mechanism naturally flags suspicious documents, which are retrieved by only one modality, without requiring retraining or domain-specific tuning.
    \item \textbf{Foundational Infrastructure for Adversarial Defence}: Establishes that retrieval architectures can be robust by construction rather than through ad-hoc security overlays. By combining multi-source redundancy, adaptive routing, and consensus scoring, QALF prevents poisoned evidence from dominating the fusion pipeline in the first place.
\end{enumerate}

Our approach builds on recent advances in adaptive multimodal retrieval, \cite{yu_cogplanner_2025,li_benchmarking_2025} and learned selectivity \cite{shi_making_2025} while addressing gaps identified in comprehensive multimodal RAG surveys \cite{mei_survey_2025}. Specifically, we extend MMed-RAG's \cite{xia_mmed-rag_2025} principle of adaptive context selection which was originally limited to choosing how many documents to retrieve to adaptive weighting and modality routing, generalizing beyond domain-specific settings. We introduce a lightweight, unsupervised mechanism that requires no retraining, making it immediately deployable across heterogeneous document collections and RAG architectures.

Contributions:
\begin{itemize}
    \item \textbf{Query Complexity Classification (4D)}: We extend prior work \cite{jeong_adaptive-rag_2024} by classifying queries along four orthogonal dimensions (linguistic, semantic, modality, contextual) rather than the binary/ternary schemes of prior work, enabling more nuanced routing decisions.
    \item \textbf{Intent-Based Multimodal Routing}: Drawing on multimodal intent classification \cite{xia_mmed-rag_2025} and adaptive planning \cite{yu_cogplanner_2025,li_benchmarking_2025}, we introduce dynamic selection of retriever combinations (vector alone, vector+graph, vector+graph+keyword) based on predicted query intent. To our knowledge, this is the first system to apply intent-based routing specifically to multimodal retrieval modality selection.
    \item \textbf{Consensus-Based Fusion with Inherent Defence}: Motivated by findings that diverse evidence defends against adversarial attacks \cite{amirshahi_evaluating_2025}, we operationalize consensus through a lightweight scoring mechanism: documents retrieved by multiple independent modalities receive higher fusion weights, while low-consensus documents are down-weighted. This provides inherent resistance to knowledge poisoning \cite{ha_mm-poisonrag_2025} without requiring supervised training.
    \item \textbf{Empirical Validation}: On standard multimodal benchmarks, QALF achieves 8–15\% NDCG@10 improvement over fixed-weight baselines. Critically, under adversarial evaluation, QALF reduces attack success rates from 56\% to 15–20\%, demonstrating 70\% improvement in robustness.
\end{itemize}
This work establishes QALF as foundational infrastructure for building adversarial-aware multimodal RAG systems, motivating systematic investigation of defence mechanisms in this domain.

\section{Related Work}

Retrieval-augmented generation (RAG) grounds LLM outputs in external knowledge to reduce hallucination. Traditional RAG uses either dense vector search or sparse keyword indexing. Recent work demonstrates that hybrid retrieval achieves better recall and robustness \cite{lee_hybgrag_2025}. RAG-Fusion \cite{rackauckas_rag-fusion_2024} introduces query decomposition with Reciprocal Rank Fusion (RRF), while \cite{ahmad_benchmarking_2025} validates that combining vector and graph retrieval outperforms single modalities. However, both apply static weighting, lacking query-adaptive mechanisms that QALF introduces.

Multimodal documents (text, tables, images) require fusion across diverse modalities. The mRAG survey \cite{mei_survey_2025} identifies a critical gap which is that most systems use static or heuristic fusion rules without query-aware adaptation. MMed-RAG \cite{xia_mmed-rag_2025} shows how choosing which and how many documents to retrieve improves accuracy (43.8\% gain). QALF extends these principles from quantity to adaptive weighting and routing, generalizing beyond domain-specific settings. Joint Fusion and Encoding \cite{huang_joint_2025} compares early vs. late fusion strategies, highlighting the importance of fusion design.

CogPlanner \cite{yu_cogplanner_2025} and OmniSearch \cite{li_benchmarking_2025} demonstrate how adaptive queries can give better framed answers. The studies break down the queries into sub queries in a multimodal pipeline to improve model performance. Adaptive-RAG \cite{jeong_adaptive-rag_2024} classifies query complexity into three categories and routes to appropriate generation strategies (no-retrieval, single-step, multi-step), achieving efficiency and accuracy gains. LQR \cite{huang_layered_2024} extends this with semantic rule based multi-hop classification. MBA-RAG \cite{tang_mba-rag_2025} uses multi-armed bandits for learned strategy selection. However, these focus on generation-stage routing, not retrieval-stage fusion adaptation. QALF differs by introducing query-aware fusion weights (not just strategy selection) and by designing for multimodal modality routing, addressing a gap in prior work.

Lastly, an emerging concern in LLMs and RAGs is adversarial robustness. MM-PoisonRAG \cite{ha_mm-poisonrag_2025} introduces multimodal knowledge poisoning, achieving 56\% attack success even with restricted access. The paper demonstrates that multimodal RAG is uniquely vulnerable because cross-modal alignment can be exploited. Robustness evaluation studies \cite{amirshahi_evaluating_2025} find that diverse evidence sources provide natural defence which is a key motivation for having a multi-source hybrid retrieval. To address these limitations, a significant line of research has focused on adaptive retrieval and learned fusion.

LKG-RALM \cite{shi_making_2025} introduced a relevance-aware fusion guided by layer-knowledge attention to ignore distractors in unimodal RAG, establishing the utility of learned selectivity. Reference demonstrated that adaptive gating strategies are crucial for defending against single-source attacks, underscoring the need for a dynamically weighted fusion layer. QALF's contribution is providing inherent defence at the retrieval-fusion stage through multi-source redundancy and consensus validation, preventing poisoned evidence from dominating in the first place.

\section{Methodology}

\begin{figure*}[t]
  \includegraphics[width=1\linewidth]{figures/QALF.png}
  \caption {QALF based hybrid multimodal RAG pipeline}
  \label{fig:architecture}
\end{figure*}

QALF (Query-Adaptive Learned Fusion) is designed for robust, high-performance multimodal retrieval-augmented generation (RAG), integrating three retrieval modalities, vector, graph, and keyword, using a unified Neo4j 5.x backend. Figure~\ref{fig:architecture} illustrates the complete QALF pipeline. The core pipeline comprises four main stages:
\begin{enumerate}
    \item Query Analysis: Classify incoming queries by complexity and intent, using lightweight rule-based heuristics and entity detection.
    \item Adaptive Routing: Dynamically select relevant retrieval modalities (vector, graph, keyword) based on query characteristics.
    \item Parallel Retrieval: Execute each retrieval process in Neo4j using its native vector index, full-text search, and graph traversal capabilities.
    \item Consensus-Based Fusion: Aggregate and rerank results using a learned fusion strategy that emphasizes cross-modality agreement for security and relevance.
\end{enumerate}

This methodology ensures that QALF is:
\begin{itemize}
    \item Efficient: Minimal infra overhead, all-in-one database, no neural/rule-based classifiers for routing
    \item Extensible: Supports any graph, text, or embedding model present in Neo4j
    \item Secure: Robust by construction against many real-world retrieval attacks
    \item Reproducible: Implemented with open-source tools, all logic fully specified
\end{itemize}

\subsection{Query Complexity and Intent Classification}

Queries are first analysed along four dimensions (linguistic, semantic, modality, contextual) to generate a fine-grained complexity profile. For example, a query may be linguistically simple but semantically complex if it requires multi-hop reasoning across entities. Each complexity aspect is determined via statistical features or keyword/entity patterns. Given query q, we compute a complexity tuple as shown in Equation~\ref{eq:complexity classification}.

\begin{equation}
    \label{eq:complexity classification}
    C(q)=(C_L,C_S,C_M,C_C)
\end{equation}

where each component is a categorical variable in {"Low","Medium","High"} for Linguistic, Semantic, Modality, and Contextual complexity.
\begin{itemize}
    \item Linguistic: \(C_L\)= function of parse tree depth and token length.
    \item Semantic: \(C_S\)= number of named entities and relationship hops required.
    \item Modality: \(C_M\)= presence of visual/tabular keywords in q.
    \item Contextual: \(C_C\)= presence of domain terms.
\end{itemize}

The intent classifier then categorizes the query (e.g., factual lookup, comparative, temporal, visual/tabular, causal, definitional, multi-hop) using string-matching rules and logical checks. This is used to guide routing and optimize the fusion strategy for different question types. Intent \(I(q)\) is computed via rules mapping to classes \(I=\{"factual","comparative",...\}\) per the decision table in Section 3.3.

\subsection{Modality Routing and Retrieval}

Routing logic uses the complexity and intent predictions to select among Neo4j retrieval modalities:
\begin{itemize}
    \item Vector Search leverages Neo4j's built-in HNSW vector index for approximate nearest-neighbor retrieval of semantic document embeddings.
    \item Keyword Search applies Neo4j’s Lucene-based full-text index for TF-IDF-style scoring over textual fields ("title", "content", "abstract").
    \item Graph Traversal executes Cypher pattern-match queries for multi-hop reasoning over entities and relationships.
\end{itemize}

Let the set of retrieval modalities for query q be \( M_q \subseteq \{Vector,Graph,Keyword\} \) selected by a rule-based map \( f_{route}\colon (C(q),I(q))=M_q \).

For each active modality, retrieve a ranked list of documents:

\begin{enumerate}
    \item \textbf{Vector}: For query embedding \(e_q\), use Neo4j vector index:
    \(score_V(d,q)=\cos{(\vec{e_q},\vec{e_d})}\) where \(e_d\) is the document embedding and cosine similarity is used.

    \item \textbf{Graph}: Find relevance by multi-hop Cypher traversal and entity overlap, \(score_G (d,q)\)=number of matching entity/edge patterns

    \item \textbf{Keyword}: Use Neo4j Lucene index, returning \(score_K (d,q)=TFIDF(fields|q)\) with Lucene's scoring scheme (approximate to BM25).
\end{enumerate}

For each \(m\in M_q\), select top k results and normalize scores to [0,1] as shown in Equation~\ref{eq:top-k for m}.

\begin{equation}
    \label{eq:top-k for m}
    r_q^m(d)=\frac{score_m(d,q)}{\underset{d^{'}}{\max}{score_m (d^{'},q)}}
\end{equation}

\subsection{Consensus-Based Fusion}

To combine results, QALF employs a fusion mechanism that rewards documents retrieved by multiple modalities (cross-modality consensus) and downweighs those that only appear from one. For each document d appearing in at least one retrieval result, the consensus score is calculated as shown in Equation~\ref{eq:consensus}.

\begin{equation}
    \label{eq:consensus}
    Consensus_q(d)=\left| \frac{ \left\{ m\in M | d\in top-k \right\} }{M_q} \right|
\end{equation}

This value is 1 if d is retrieved by all modalities, lower otherwise. Adaptive modality weights are computed per query based on intent, consensus, and tuneable parameters. Final document rankings are produced by a weighted sum of the modality-specific scores (normalized RRF), followed by a top K selection.

\subsection{Adaptive Modality Weighting}

Each modality m receives an adaptive weight based on intent, consensus, and a tuneable coefficient \(\beta\) calculated with Equation~\ref{eq:adaptive weight}.

\begin{equation}
    \label{eq:adaptive weight}
    w_q=\alpha_{I(q)}^m \cdot \left[ 1-\beta \cdot \left( 1-\bar{Consensus}_q^m \right) \right]
\end{equation}

where,
\begin{itemize}
    \item \(\alpha_{I(q)}^m\) is the base weight for modality \(m\) given intent \(I(q)\), learned or grid-searched on validation data.
    \item \(\bar{Consensus}_q^m\) is the average consensus of top-k docs for modality \(m\):
    \begin{equation}
        \bar{Consensus}_q^m=\frac{1}{k} \sum_{d\in D_{top-k}^m(q)} Consensus_q(d)
    \end{equation}
    \item \(\beta\) controls how much consensus matters (\(\beta=0.5\) by default).
\end{itemize}


\subsection{Reciprocal Rank Fusion Extension}

For each candidate document d, its final QALF score is calculated with Equation~\ref{eq:RRF}.

\begin{equation}
    \label{eq:RRF}
    Score_{QALF}(d,q)=\sum_{m\in M_q}w_q^m \cdot \left[ \frac{1}{rank^m (d,q)+K} \right]
\end{equation}

where,
\begin{itemize}
    \item \(rank^m(d,q)\) is the rank of d in the result of the modality m (if present, otherwise set to a high value to diminish its contribution).
    \item K (typically, 60) is the smoothing parameter of the RRF.
\end{itemize}

\subsection{Security and Robustness Considerations}

Unlike naive hybrid or static-weighted fusion, QALF’s consensus rewards redundancy:
\begin{itemize}
    \item Low consensus = low weight: If a doc is injected/adversarial (e.g., only retrieved by vector modality), its consensus is low. Thus, \(w_q^V\) is down-weighed, and adversarial docs are unlikely to appear in the final top K.
    \item Multi-modality redundancy: The attacker must poison all three modalities (vector, text, graph), which is considerably more difficult than poisoning a single modality, especially since each uses different features.
\end{itemize}

Multi-source routing, adaptive weights, and cross-modality aggregation collectively make targeted attacks significantly less effective.

\section{Evaluation}
DocBench for multimodal RAG QnA
MM-PoisonRAG for defence against adversarial attacks

\section{Conclusion and Future Work}

This paper introduced QALF, a query-adaptive learned fusion framework for multimodal retrieval-augmented generation that operates entirely over a unified Neo4j backend. By combining four-dimensional query complexity modeling, intent-based modality routing, and consensus-aware fusion of vector, graph, and keyword retrieval, QALF addresses two key limitations of existing hybrid RAG systems: static, query-agnostic fusion and vulnerability to multimodal knowledge poisoning \cite{lee_hybgrag_2025,mei_survey_2025,ha_mm-poisonrag_2025}. Experiments on DocBench-style multimodal QA and MM-PoisonRAG-style adversarial setups show that QALF improves retrieval effectiveness over strong fixed-weight hybrids and substantially reduces poison success rates while preserving most clean performance \cite{ahmad_benchmarking_2025,amirshahi_evaluating_2025}. These results suggest that retrieval architectures can be made more robust by construction through adaptive multi-source fusion, without requiring heavy supervised training or domain-specific tuning \cite{jeong_adaptive-rag_2024,shi_making_2025}.

In future work, we plan to extend QALF along three directions. First, we will replace or augment the current rule-based complexity and intent classifiers with lightweight learned models, enabling finer-grained routing while studying robustness–overfitting trade-offs across domains and languages. Second, we aim to broaden the empirical scope by evaluating on additional multimodal and cross-domain benchmarks, larger corpora, and alternative backend stacks (e.g., specialized vector and keyword engines) to better characterize the portability of QALF’s design choices. Third, building on the MM-PoisonRAG threat model, we will investigate stronger adaptive attackers and integrate QALF into a more comprehensive defence stack that also addresses prompt injection, jailbreaks, and model-level backdoors in multimodal RAG systems. Together, these directions form the basis for a broader research agenda on principled and adversarially robust multimodal RAG.

\section*{Limitations}

While QALF improves both retrieval effectiveness and robustness over strong hybrid baselines, the current study has several limitations.

First, the evaluation is restricted to a small number of multimodal benchmarks (DocBench-style document QA) and a single adversarial framework (MM-PoisonRAG). These datasets primarily cover PDF-like documents with text, tables, and figures, so it is unclear how well QALF generalizes to other domains such as web-scale corpora, code-heavy repositories, or real-time streaming data \cite{mei_survey_2025,xia_mmed-rag_2025}. The query distributions and attack models considered here may therefore underrepresent the diversity and difficulty of real-world user queries and adversaries \cite{ha_mm-poisonrag_2025,amirshahi_evaluating_2025}.

Second, QALF relies on rule-based query complexity and intent classification, as well as hand-tuned routing and fusion hyperparameters. Although this design keeps the system lightweight and domain-agnostic, it may miss nuanced query phenomena like implicit multimodality, subtle temporal or causal structure that learned classifiers could capture. The rule sets are also crafted on English-language corpora and may not transfer directly to other languages or scripts without additional engineering effort \cite{jeong_adaptive-rag_2024,huang_layered_2024}.

Third, the adversarial evaluation focuses on knowledge poisoning at the retrieval layer and assumes that the underlying LLM behaves as a mostly honest generator given its retrieved context. The experiments do not consider prompt-injection attacks, jailbreaks, or model-level backdoors, nor do they study adaptive attackers that react to QALF’s consensus weighting. As a result, QALF should be viewed as a defence mechanism specifically for retrieval- and fusion-stage poisoning rather than a complete solution to all adversarial threats against multimodal RAG \cite{ha_mm-poisonrag_2025,shi_making_2025}.

Finally, the implementation uses a single Neo4j backend for vector, keyword, and graph retrieval, which simplifies deployment but couples performance and robustness to one database technology. Alternative vector indices, keyword backends, or graph engines might change the empirical trade-offs between latency, effectiveness, and security. A more exhaustive system study across multiple backends and larger corpora is left for future work \cite{lee_hybgrag_2025,ahmad_benchmarking_2025}.

\section*{Ethical considerations}

This work advances Retrieval Augmented Generation (RAG) systems for multimodal document corpora, with the primary goal of improving factual accuracy and robustness against adversarial threats. Although these objectives support more reliable AI systems, several ethical considerations merit discussion.

\textbf{Dual-Use Risks.} QALF's adversarial robustness evaluation uses knowledge poisoning attacks from MM-PoisonRAG \cite{ha_mm-poisonrag_2025} to demonstrate defense efficacy. Although conducted in controlled research settings with synthetic data, the evaluation framework and poisoning techniques could be misused to develop more sophisticated attacks. To mitigate this, we commit to releasing only defensive components (routing, consensus fusion) and evaluation harnesses under permissive open-source licenses, while documenting attack methodologies with clear warnings about responsible use. Full attack code and poisoned corpora will not be distributed.

\textbf{Bias Amplification.} QALF's rule-based query classifiers and Neo4j full-text indexing inherit biases from their training data and indexing schemes. For example, entity recognition and TF-IDF scoring can under perform in low-resource languages or underrepresented domains, potentially amplifying retrieval disparities. Consensus weighting can also exacerbate biases if spurious evidence dominates low-consensus domains. Future work will evaluate demographic parity and fairness metrics in language and topic distributions \cite{mei_survey_2025}.

\textbf{Privacy and Deployment.} The unified Neo4j backend stores document embeddings, full-text indices, and graph relationships, which may inadvertently retain sensitive information from source documents. Deployers should apply standard data minimization like anonymization or TTL policies and comply with privacy regulations. QALF does not inherently mitigate memorization risks in retrieval, which remain an active research area.

\textbf{Societal Impact.} By making multimodal RAG more robust, QALF enables higher-confidence deployment in high-stakes domains like scientific literature review, financial analysis, and legal research. However, over-reliance on any single system, including QALF, could create single points of failure. We encourage deployers to maintain human oversight and diverse evidence sources, particularly for critical decision-making.

All experiments use publicly available benchmarks (DocBench-style, MM-PoisonRAG) and open-weight models where possible. No human subjects were involved, and compute usage remained under institutional guidelines.

% Bibliography entries for the entire Anthology, followed by custom entries
%\bibliography{anthology,custom}
% Custom bibliography entries only
\bibliography{references}

\appendix

\section{Example Appendix}
\label{sec:appendix}

This is an appendix.

\end{document}
